% GENERATED FILE. Edit canonical lessons, then run npm run generate:books.
% canonical-source-hash: fnv1a64:58de0287c80b232b
% canonical-lessons: AR-C552-tila, AR-C552-rammama, AR-C552-zayyana, AR-C552-jaffafa, AR-C552-farasha, AR-R552-review-552-552, AR-R552-review-final-552

\chapter{Doing up the house}
\label{ch:ar-doing-up-the-house}
\hlchaptermodality{552}

\begin{chapteropening}
\textbf{By the end of this chapter:} \emph{I can talk about a coat of paint, and say that someone renovated, decorated, furnished or dried something.}

Complete the last lesson of chapter 552: I can talk about a coat of paint, and say that someone renovated, decorated, furnished or dried something.
\end{chapteropening}

\section[\texorpdfstring{\ar{طلاء} (ṭilāʾ)}{طلاء (ṭilāʾ)}]{\ar{طلاء} (ṭilāʾ) — a coat of paint}

\label{lesson:AR-C552-tila}



\begin{quote}
Before the new one: say the Arabic for a mop, then the Arabic for a couch.
\end{quote}

\subsection*{You'll want to know: \ar{طلاء}}
\textbf{\ar{طلاء}} — \emph{ṭilāʾ} — \textquotedblleft{}a coat of paint\textquotedblright{}.

\begin{quote}
\ar{جف} \ar{الطلاء} \ar{الأزرق} \ar{على} \ar{الباب}.
\end{quote}

\begin{quote}
\emph{jaffa aṭ-ṭilāʾ al-azraq ʿalā al-bāb.}
\end{quote}

\begin{quote}
The blue paint dried on the door.
\end{quote}

\begin{cousinweb}
\textbf{\ar{طلاء}} (\emph{ṭilāʾ}) is built on the native Arabic root \emph{ṭ-l-y}, \textquotedblleft{}to smear, to coat\textquotedblright{} (Lane). Ṭilāʾ is what is smeared on.
\end{cousinweb}

\begin{grammarlens}[title={What you've built}]
One more word for describing what you lived and felt.
\end{grammarlens}

\subsection*{Your turn}
\begin{itemize}
\raggedright
  \item \emph{Say it:} \emph{ṭilāʾ}
  \item \emph{Say it:} \emph{ṭilāʾ}, once more
  \item \emph{Say it:} \emph{jaffa aṭ-ṭilāʾ al-azraq ʿalā al-bāb.} — The blue paint dried on the door.
  \item \emph{Recall:} say the Arabic for a mop, then the Arabic for a couch, then say \emph{ṭilāʾ} again
\end{itemize}

\begin{tcolorbox}[breakable,colback=teal!4,colframe=teal!35!black,title={Before you move on}]
What does \ar{طلاء} mean? (\textquotedblleft{}A coat of paint\textquotedblright{}.) Say it once more.
\end{tcolorbox}
\section[\texorpdfstring{\ar{رمم} (rammama)}{رمم (rammama)}]{\ar{رمم} (rammama) — to renovate, to restore}

\label{lesson:AR-C552-rammama}



\begin{quote}
Before the new one: say the Arabic for a couch, then the Arabic for a coat of paint.
\end{quote}

\subsection*{You'll want to know: \ar{رمم}}
\textbf{\ar{رمم}} — \emph{rammama} — \textquotedblleft{}to renovate, to restore\textquotedblright{}.

\begin{quote}
\ar{رمم} \ar{أبي} \ar{البيت} \ar{القديم}.
\end{quote}

\begin{quote}
\emph{rammama abī al-bayt al-qadīm.}
\end{quote}

\begin{quote}
My father renovated the old house.
\end{quote}

Present: \ar{يرمم}, \emph{yurammim}.

\begin{cousinweb}
\textbf{\ar{رمم}} (\emph{rammama}) is built on the native Arabic root \emph{r-m-m}, \textquotedblleft{}to repair, to mend\textquotedblright{} (Lane). Form II rammama is to put an old thing right.
\end{cousinweb}

\begin{grammarlens}[title={What you've built}]
One more word for describing what you lived and felt.
\end{grammarlens}

\subsection*{Your turn}
\begin{itemize}
\raggedright
  \item \emph{Say it:} \emph{rammama}
  \item \emph{Say it:} \emph{rammama}, once more
  \item \emph{Say it:} \emph{rammama abī al-bayt al-qadīm.} — My father renovated the old house.
  \item \emph{Recall:} say the Arabic for a couch, then the Arabic for a coat of paint, then say \emph{rammama} again
\end{itemize}

\begin{tcolorbox}[breakable,colback=teal!4,colframe=teal!35!black,title={Before you move on}]
What does \ar{رمم} mean? (\textquotedblleft{}To renovate, to restore\textquotedblright{}.) Say it once more.
\end{tcolorbox}
\section[\texorpdfstring{\ar{زين} (zayyana)}{زين (zayyana)}]{\ar{زين} (zayyana) — to decorate}

\label{lesson:AR-C552-zayyana}



\begin{quote}
Before the new one: say the Arabic for a coat of paint, then the Arabic for to renovate.
\end{quote}

\subsection*{You'll want to know: \ar{زين}}
\textbf{\ar{زين}} — \emph{zayyana} — \textquotedblleft{}to decorate\textquotedblright{}.

\begin{quote}
\ar{زين} \ar{الأطفال} \ar{الغرفة} \ar{للعيد}.
\end{quote}

\begin{quote}
\emph{zayyana al-aṭfāl al-ghurfa li-l-ʿīd.}
\end{quote}

\begin{quote}
The children decorated the room for the feast.
\end{quote}

Present: \ar{يزين}, \emph{yuzayyin}.

\begin{cousinweb}
\textbf{\ar{زين}} (\emph{zayyana}) is built on the native Arabic root \emph{z-y-n}, \textquotedblleft{}to adorn, to beautify\textquotedblright{} (Lane). Zīna, decorations, which you know, is from the same root.
\end{cousinweb}

\begin{grammarlens}[title={What you've built}]
One more word for describing what you lived and felt.
\end{grammarlens}

\subsection*{Your turn}
\begin{itemize}
\raggedright
  \item \emph{Say it:} \emph{zayyana}
  \item \emph{Say it:} \emph{zayyana}, once more
  \item \emph{Say it:} \emph{zayyana al-aṭfāl al-ghurfa li-l-ʿīd.} — The children decorated the room for the feast.
  \item \emph{Recall:} say the Arabic for a coat of paint, then the Arabic for to renovate, then say \emph{zayyana} again
\end{itemize}

\begin{tcolorbox}[breakable,colback=teal!4,colframe=teal!35!black,title={Before you move on}]
What does \ar{زين} mean? (\textquotedblleft{}To decorate\textquotedblright{}.) Say it once more.
\end{tcolorbox}
\section[\texorpdfstring{\ar{جفف} (jaffafa)}{جفف (jaffafa)}]{\ar{جفف} (jaffafa) — to dry off}

\label{lesson:AR-C552-jaffafa}



\begin{quote}
Before the new one: say the Arabic for to renovate, then the Arabic for to decorate.
\end{quote}

\subsection*{You'll want to know: \ar{جفف}}
\textbf{\ar{جفف}} — \emph{jaffafa} — \textquotedblleft{}to dry off\textquotedblright{}.

\begin{quote}
\ar{جففت} \ar{أمي} \ar{الأطباق} \ar{بمنشفة}.
\end{quote}

\begin{quote}
\emph{jaffafat ummī al-aṭbāq bi-minshafa.}
\end{quote}

\begin{quote}
My mother dried the plates with a towel.
\end{quote}

Present: \ar{يجفف}, \emph{yujaffif}.

\begin{cousinweb}
\textbf{\ar{جفف}} (\emph{jaffafa}) is built on the native Arabic root \emph{j-f-f}, \textquotedblleft{}to be dry\textquotedblright{} (Lane). Jāff, dry, which you know, is from the same root; Form II jaffafa is to make dry.
\end{cousinweb}

\begin{grammarlens}[title={What you've built}]
One more word for describing what you lived and felt.
\end{grammarlens}

\subsection*{Your turn}
\begin{itemize}
\raggedright
  \item \emph{Say it:} \emph{jaffafa}
  \item \emph{Say it:} \emph{jaffafa}, once more
  \item \emph{Say it:} \emph{jaffafat ummī al-aṭbāq bi-minshafa.} — My mother dried the plates with a towel.
  \item \emph{Recall:} say the Arabic for to renovate, then the Arabic for to decorate, then say \emph{jaffafa} again
\end{itemize}

\begin{tcolorbox}[breakable,colback=teal!4,colframe=teal!35!black,title={Before you move on}]
What does \ar{جفف} mean? (\textquotedblleft{}To dry off\textquotedblright{}.) Say it once more.
\end{tcolorbox}
\section[\texorpdfstring{\ar{فرش} (farasha)}{فرش (farasha)}]{\ar{فرش} (farasha) — to furnish; to spread}

\label{lesson:AR-C552-farasha}



\begin{quote}
Before the new one: say the Arabic for to decorate, then the Arabic for to dry off.
\end{quote}

\subsection*{You'll want to know: \ar{فرش}}
\textbf{\ar{فرش}} — \emph{farasha} — \textquotedblleft{}to furnish; to spread\textquotedblright{}.

\begin{quote}
\ar{فرشنا} \ar{الشقة} \ar{الجديدة}.
\end{quote}

\begin{quote}
\emph{farashnā ash-shaqqa al-jadīda.}
\end{quote}

\begin{quote}
We furnished the new flat.
\end{quote}

Present: \ar{يفرش}, \emph{yafrush}.

\begin{cousinweb}
\textbf{\ar{فرش}} (\emph{farasha}) is built on the native Arabic root \emph{f-r-sh}, \textquotedblleft{}to spread out\textquotedblright{} (Lane). Firāsh, a mattress, which you know, is from the same root: a thing spread out.
\end{cousinweb}

\begin{grammarlens}[title={What you've built}]
One more word for describing what you lived and felt.
\end{grammarlens}

\subsection*{Your turn}
\begin{itemize}
\raggedright
  \item \emph{Say it:} \emph{farasha}
  \item \emph{Say it:} \emph{farasha}, once more
  \item \emph{Say it:} \emph{farashnā ash-shaqqa al-jadīda.} — We furnished the new flat.
  \item \emph{Recall:} say the Arabic for to decorate, then the Arabic for to dry off, then say \emph{farasha} again
\end{itemize}

\begin{tcolorbox}[breakable,colback=teal!4,colframe=teal!35!black,title={Before you move on}]
What does \ar{فرش} mean? (\textquotedblleft{}To furnish; to spread\textquotedblright{}.) Say it once more.
\end{tcolorbox}
\section[(dialogue)]{The house — a last review}

\label{lesson:AR-R552-review-552-552}



\begin{quote}
Say the Arabic for a coat of paint and to renovate.
\end{quote}

\subsection*{Your turn}
\begin{itemize}
\raggedright
  \item \emph{Say it:} a coat of paint, to renovate, to decorate, to dry off, to furnish
\end{itemize}

\begin{tcolorbox}[breakable,colback=teal!4,colframe=teal!35!black,title={Before you move on}]
A coat of paint? (\textbf{\ar{طلاء}}, \emph{ṭilāʾ}.) To renovate? (\textbf{\ar{رمم}}, \emph{rammama}.) To decorate? (\textbf{\ar{زين}}, \emph{zayyana}.) To dry off? (\textbf{\ar{جفف}}, \emph{jaffafa}.) To furnish? (\textbf{\ar{فرش}}, \emph{farasha}.)
\end{tcolorbox}
\section[(dialogue)]{Doing up the house — once more}

\label{lesson:AR-R552-review-final-552}



\begin{quote}
Say the Arabic for a coat of paint and to renovate.
\end{quote}

\subsection*{Your turn}
\begin{itemize}
\raggedright
  \item \emph{Say it:} a coat of paint, to renovate, to decorate, to dry off, to furnish
\end{itemize}

\begin{tcolorbox}[breakable,colback=teal!4,colframe=teal!35!black,title={Before you move on}]
A coat of paint? (\textbf{\ar{طلاء}}, \emph{ṭilāʾ}.) To renovate? (\textbf{\ar{رمم}}, \emph{rammama}.) To decorate? (\textbf{\ar{زين}}, \emph{zayyana}.) To dry off? (\textbf{\ar{جفف}}, \emph{jaffafa}.) To furnish? (\textbf{\ar{فرش}}, \emph{farasha}.)
\end{tcolorbox}
